\section{Bedienungsanleitung}
\label{sec:bedienungsanleitung}
Die folgende Kurzanleitung fasst die Bedienung der Temperaturanzeige für den Bediener zusammen (vgl. Kapitel~\ref{sec:bedienkonzept}).

\begin{enumerate}[itemsep=2pt]
    \item \textbf{Einschalten:} Nach dem SPS-Start fährt der Zeiger automatisch auf die Nullposition (beide LEDs blinken schnell). Danach wird die aktuelle Temperatur angezeigt (grüne LED blinkt langsam).
    \item \textbf{Ablesen:} Skalenwert = Temperatur in 1/10\,\textdegree C (z.\,B. 250 = 25,0\,\textdegree C). Der Zeiger wird jede Sekunde automatisch nachgeführt.
    \item \textbf{Zeiger korrigieren:} Grünen Taster länger als 300\,ms halten, um den Zeiger nach oben zu fahren, roten Taster für nach unten. Beim Loslassen wird die Korrektur gespeichert. Am Skalenende stoppt der Zeiger automatisch.
    \item \textbf{Neu starten:} Beide Taster 3\,s gleichzeitig halten. Die Korrektur wird gelöscht, der Zeiger fährt auf Null und danach auf die aktuelle Temperatur (beide LEDs blinken schnell).
    \item \textbf{Beide LEDs blinken langsam:} Temperatur liegt außerhalb des Anzeigebereichs (0--36\,\textdegree C), kein Eingriff nötig.
    \item \textbf{Nur rote LED blinkt:} Fehler. Prüfen, ob der Zeiger frei beweglich ist und der Sensor angeschlossen ist. Dann beide Taster 3\,s halten. Tritt der Fehler erneut auf, Instandhaltung informieren.
\end{enumerate}
